\chapter{Sistemas y memoria}
\label{cap:sistemas}

\begin{intuicion}
Los algoritmos de matching que vamos a estudiar son inteligentes,
pero la mitad del rendimiento real del motor sale de decisiones de
\emph{layout} de memoria: \texttt{CompactRelation} pesa 32~bytes en
vez de 200 porque cada arista que cabe en la \emph{cache} es una que
no espera a la RAM; el \emph{streaming backend} usa \emph{mmap}
porque transferir 28~GB del disco al heap es inviable; el \emph{lazy
build} de NLF y \emph{khop\_reach} usa \texttt{OnceLock} porque la
inicialización doble bajo concurrencia es un bug clásico. Este
capítulo da los fundamentos de sistemas que esas decisiones invocan.
Si ya sabés qué es una \emph{cache line}, qué garantiza un
\emph{atomic}, y por qué un \emph{branch mispredict} cuesta ciclos,
podés saltearlo. Si no, leélo: el resto del libro lo asume.
\end{intuicion}

\section{La jerarquía de memoria}

Una operación aritmética en un núcleo de CPU moderno cuesta
\(\sim 1\) ciclo. Acceder a la RAM cuesta \(\sim 200\) ciclos. Esa
diferencia ---\emph{un orden de magnitud y medio}--- es lo que
manda toda la conversación sobre \emph{layout} de datos. Para
amortiguar, las CPUs tienen una jerarquía de \emph{caches}:

\begin{center}
\renewcommand{\arraystretch}{1.3}
\begin{tabular}{lrrr}
\toprule
\textbf{Nivel} & \textbf{Tamaño típico} & \textbf{Latencia (ciclos)} & \textbf{Latencia (ns)} \\
\midrule
Registros        & 100 B             & 0           & 0 \\
L1 cache         & 32--64 KB         & 4--5        & 1 \\
L2 cache         & 256 KB--1 MB      & 10--12      & 3 \\
L3 cache         & 8--32 MB          & 30--40      & 12 \\
RAM principal    & 16--128 GB        & 200--300    & 100 \\
SSD NVMe         & 1--2 TB           & \(\sim 10^5\) & 30--80 \(\mu\)s \\
HDD              & 8--16 TB          & \(\sim 10^7\) & 4--10 ms \\
\bottomrule
\end{tabular}
\end{center}

\noindent
La regla de oro: si tus datos caben en L2, los algoritmos
``inteligentes'' rinden lo que prometen. Si no caben, la latencia de
RAM domina y la ``inteligencia'' algorítmica importa menos que la
\emph{locality} de acceso.

\subsection{Cache lines}

La CPU no transfiere bytes individuales entre RAM y cache: transfiere
\emph{cache lines} de 64~bytes (en x86\_64 e ARM modernos). Acceder a
un solo byte trae los 63 contiguos ``gratis''. Eso tiene dos
implicaciones:

\begin{enumerate}
    \item \textbf{Datos contiguos son rápidos.} Iterar un
    \texttt{Vec<u32>} de 1000 elementos toca \(1000 \cdot 4 / 64 =
    62.5\) cache lines. El siguiente elemento ya está cargado cuando
    el anterior se procesa (la \emph{prefetch} unit lo trae).
    \item \textbf{Datos esparcidos son lentos.} Iterar un
    \texttt{HashMap<u32, u32>} de 1000 elementos toca, en el peor
    caso, 1000 cache lines distintas (cada \emph{bucket} en una
    región distinta de RAM). \(60\times\) más cache misses, mismo
    trabajo lógico.
\end{enumerate}

\begin{enelcodigo}
Esta es la justificación de fondo de por qué
\texttt{CompactRelation} es 32 bytes (\emph{half} de una cache line:
dos aristas por línea). En el DFS hot loop, iterar
\texttt{neighbors\_of(v)} carga 32 aristas por cada 16 cache lines
---versus el \texttt{Relation} público de \(\sim 200\) bytes que
cargaría 32 aristas en \(\sim 100\) cache lines.
\filepath{core/graph-engine/src/types.rs}.
\end{enelcodigo}

\subsection{Branch prediction}

Las CPUs modernas no esperan a que un \texttt{if} se evalúe para
decidir qué instrucción cargar. Predicen el resultado y lo ejecutan
\emph{especulativamente}; si la predicción acierta, ganaron 10--20
ciclos; si falla, descartan el trabajo y empiezan de nuevo (penalty
de 15--25 ciclos en x86\_64).

\begin{itemize}
    \item Loops contables (\texttt{for i in 0..n}): predicción casi
    perfecta. Branch barato.
    \item Branches dependientes de datos
    impredecibles (\texttt{if data[i] > threshold}): mispredicts
    frecuentes si los datos son aleatorios.
    \item Hot loops con muchos \texttt{if} anidados: cada
    \texttt{if} es una oportunidad de mispredict; los \emph{loops}
    apretados con un solo branch son de oro.
\end{itemize}

\begin{nota}
Esta es la razón por la que el inner loop del DFS de GraphHunter está
diseñado como una sola pasada lineal sobre un \emph{slice} contiguo
de \texttt{CompactRelation}, con los predicados temporales reducidos
a comparaciones de \texttt{i64} y los predicados de tipo a
comparaciones de \texttt{u8}. Los branches que quedan son
predecibles por la distribución típica de los datos (la mayoría de
las aristas pasan los filtros tempranos).
\end{nota}

\subsection{SIMD}

Las CPUs x86\_64 e ARM tienen instrucciones que operan sobre 4, 8 o
16 valores a la vez (\emph{SIMD}: Single Instruction, Multiple Data).
Para problemas que se mapean naturalmente a operaciones por elemento
---\emph{intersección de conjuntos ordenados}, por ejemplo--- el
\emph{lift} versus la versión escalar es proporcional al ancho del
vector.

GraphHunter usa SIMD vía \texttt{libgraphmatch} para la
intersección \emph{multi-way} en LFTJ-K\(_4\)
(capítulo~\ref{cap:lftj}): \(N(a) \cap N(b) \cap N(c)\) sobre
adyacencias ordenadas. Con AVX2, intersecta 8 \texttt{u32}s por
ciclo; con AVX-512, 16.

\section{Memoria virtual y mmap}

\subsection{Memoria virtual: la abstracción del SO}

Los procesos no ven la RAM física directamente. Ven un \emph{address
space} virtual de 64~bits que el kernel mapea on-demand a páginas
físicas (típicamente 4~KB cada una). Beneficios:

\begin{itemize}
    \item Cada proceso cree que tiene 16 EB de memoria.
    \item El SO mueve páginas entre RAM y \emph{swap} sin que el
    proceso se entere.
    \item Las regiones de archivos pueden \emph{mapearse} al address
    space (mmap): el archivo en disco aparece como bytes accesibles
    con punteros.
\end{itemize}

\subsection{mmap}

\begin{definition}[Memory-mapped file]
Una región mmap es una asociación entre un rango de páginas
virtuales y un rango de bytes en un archivo. Acceder a una página no
residente dispara un \emph{page fault}: el kernel lee los 4~KB del
archivo, los pone en una página física libre, y reanuda la
ejecución.
\end{definition}

Tres garantías importantes:

\begin{enumerate}
    \item \textbf{Lazy loading.} El archivo no se lee hasta que el
    proceso accede a una página. Mapear un archivo de 28~GB ocupa
    \emph{nada} de RAM; sólo registra un descriptor.
    \item \textbf{Cache compartida con el page cache del SO.} Las
    páginas leídas viven en el \emph{page cache}; otros procesos que
    mapean el mismo archivo las comparten.
    \item \textbf{Backpressure automático.} Si el sistema necesita
    RAM, el kernel descarta páginas mmap-resident
    (sin escribir nada al disco si están \emph{clean}); el próximo
    acceso las relee. Es \emph{free swap} para datos read-only.
\end{enumerate}

\begin{advertencia}
El acceso aleatorio sobre mmap es lento si dispara muchos page
faults (cada uno cuesta \(\sim 100\,000\) ciclos). El acceso
\emph{secuencial} es rápido porque el kernel hace
\emph{readahead}: cuando ve un patrón secuencial, pre-carga las
próximas páginas en background. Por eso GraphHunter ordena las
aristas por \texttt{(source\_sid, timestamp)} antes de hacer spill:
los \emph{accesses} del matcher son secuenciales por construcción.
\end{advertencia}

\begin{enelcodigo}
\filepath{core/graph-engine/src/streaming.rs}: el \emph{streaming
backend} aplica el mismo principio. Cuando el budget global
(default $10^6$ aristas) se cruza, cada \emph{per-vertex slice}
se vuelca sortado por timestamp a un archivo temporal y la
\texttt{NeighborList} de cada vértice swap-ea a la variante
\texttt{Mmap}. Los lookups siguientes van vía puntero al archivo;
el SO decide qué páginas viven en RAM.
\end{enelcodigo}

\section{Concurrencia: atomics, locks, y wait-free}

\subsection{Lectura concurrente: el problema}

GraphHunter es multi-thread: el ingestor escribe aristas mientras
los matchers de varios transports leen. La pregunta es: ¿cómo
evitamos que un \emph{reader} vea estado a medio escribir?

Tres estrategias clásicas:

\begin{enumerate}
    \item \textbf{Mutex.} Un solo thread tiene acceso a la vez.
    Simple, correcto, pero serializa todo: si hay 100 readers, 99
    están esperando.
    \item \textbf{Read-Write lock (\texttt{RwLock}).} Múltiples
    readers concurrentes, un único writer exclusivo. Mejor para
    workloads read-heavy.
    \item \textbf{Wait-free.} Los readers nunca bloquean: ven
    siempre un \emph{snapshot consistente} aunque el writer esté
    activo. Implementación típica: el writer construye una nueva
    versión y publica el puntero atómicamente; los readers ven la
    versión vieja o la nueva, nunca un estado intermedio.
\end{enumerate}

\subsection{Operaciones atómicas}

Las CPUs ofrecen instrucciones que ejecutan
\emph{lectura-modificación-escritura} sobre una palabra de memoria
sin que otro thread pueda interrumpir. Las primitivas más comunes:

\begin{itemize}
    \item \texttt{load}/\texttt{store} atómicos: leer o escribir una
    palabra de modo que ningún reader vea ``mitad de un valor''.
    \item \texttt{compare\_and\_swap (CAS)}: ``si el valor actual es
    \(x\), reemplazalo por \(y\), todo de una''. Devuelve el valor
    viejo (o un \emph{success} flag).
    \item \texttt{fetch\_add}: incrementar atómicamente, devolver el
    valor previo. Útil para \emph{counters} compartidos.
\end{itemize}

Las atómicas son \(\sim 5\)--\(20\times\) más caras que sus
contrapartes no-atómicas, pero son \(100\times\) más baratas que un
\texttt{Mutex} bajo contención.

\subsection{Memory ordering}

\begin{advertencia}
Acá hay una sutileza profunda: una CPU moderna \emph{reordena}
operaciones de memoria para ganar performance. Sin garantías
explícitas, un thread puede ver las escrituras de otro
\emph{en orden distinto al que se ejecutaron}. Las atómicas vienen
con un parámetro de \emph{memory ordering} (\texttt{Relaxed},
\texttt{Acquire}, \texttt{Release}, \texttt{SeqCst}) que controla
qué reordenamientos están permitidos. Usar \texttt{Relaxed} cuando
hace falta \texttt{Acquire/Release} es un bug clásico que NO se
manifiesta en hardware débil (x86) pero sí en ARM.
\end{advertencia}

GraphHunter usa \texttt{OnceLock} (de la stdlib de Rust) para los
índices \emph{lazy-built}: el primer thread que llama
\texttt{ensure\_nlf()} construye el índice; los demás ven el
resultado vía un \texttt{Acquire} load del puntero. Sin \emph{double
init}, sin \emph{tornados} de constructores.

\subsection{Wait-free para el streaming backend}

El \emph{streaming backend} (capítulo~\ref{cap:streaming-mirror})
soporta \emph{wait-free reads} via una estructura
\emph{degree-adaptive}: cada vértice tiene un \texttt{NeighborList}
que es \texttt{Inline}, \texttt{Vec}, o \texttt{Btree} según su
grado. El writer publica nuevas aristas vía un puntero atómico; el
reader ve el snapshot que existía cuando entró. Sin lock.

\section{Layout de datos en Rust}

\subsection{\texttt{\#[repr(C)]} y predictibilidad}

Por defecto, el compilador de Rust puede reordenar los campos de un
\texttt{struct} para minimizar padding. Eso ahorra memoria en algunos
casos, pero rompe la asunción de que el \texttt{struct} se mapea
directamente a un buffer binario.

\texttt{\#[repr(C)]} fuerza el orden declarado y las reglas de
alineación de C: cada campo va exactamente donde el programador
quiere, con padding explícito si hace falta. Esto permite:

\begin{itemize}
    \item Mapear \texttt{Vec<CompactRelation>} directamente a un
    archivo mmap: los bytes en disco son los bytes en memoria.
    \item Pasar \texttt{CompactRelation} por FFI a C/C++ sin
    \emph{marshalling}.
    \item Razonar sobre el tamaño exacto (32 bytes, no ``alguna vez
    24 alguna vez 40 según versión del compilador'').
\end{itemize}

\begin{ejemplo}
Sin \texttt{repr(C)}:
\begin{lstlisting}[language=Rust]
struct Bad {
    flag: u8,        // 1 B
    timestamp: i64,  // 8 B
    sid: u32,        // 4 B
}
\end{lstlisting}
El compilador puede ordenar como \texttt{(timestamp, sid, flag)} para
ahorrar padding ---o no---. Tamaño: 16 B (con padding) o 24 B
(reordenado). Impredecible.

Con \texttt{repr(C)}:
\begin{lstlisting}[language=Rust]
#[repr(C)]
struct Good {
    timestamp: i64,  // offset 0,  8 B
    sid: u32,        // offset 8,  4 B
    flag: u8,        // offset 12, 1 B
    _pad: [u8; 3],   // offset 13, 3 B padding explícito
}
\end{lstlisting}
Total: 16 B exactos. Layout estable a través de versiones del
compilador. El programador eligió.
\end{ejemplo}

\subsection{\texttt{Arc<str>} y \emph{refcounting}}

\texttt{Arc<T>} (\emph{atomic reference count}) es un puntero
compartido que mantiene un contador atómico de cuántos dueños
tiene. Clonar un \texttt{Arc} no copia los datos; sólo bumpea el
contador (un \texttt{fetch\_add} atómico, \(\sim 5\) ns). Cuando el
último dueño se va, los datos se liberan.

\texttt{Arc<str>} es la versión para strings: una sola copia del
texto en heap, refcount atómico, clonado barato. La diferencia
con \texttt{String::clone()} (que copia los bytes) es del orden de
\emph{decenas a cientos de nanosegundos por clone} ---suficiente para
explicar el \(55\%\) de mejora en cap-saturated benches que motivó
el cambio de \texttt{HuntResult = Vec<String>} a
\texttt{Vec<Arc<str>>} (capítulo~\ref{cap:hunt-result}).

\section{Cómo todo esto vuelve en el resto del libro}

Los conceptos de este capítulo aparecen explícitamente en:

\begin{itemize}
    \item Capítulo~\ref{cap:compact-relation}: \texttt{CompactRelation}
    32 B, \texttt{repr(C)}, half-cache-line.
    \item Capítulo~\ref{cap:streaming-mirror}: el \emph{streaming
    backend} con auto-spill, mmap, y page cache.
    \item Capítulo~\ref{cap:hunt-result}: migración a
    \texttt{Vec<Arc<str>>}, refcount versus alloc.
    \item Capítulo~\ref{cap:streaming-mirror}: streaming backend
    wait-free, \texttt{neighbors\_snapshot}.
    \item Capítulo~\ref{cap:incremental-nlf}: \texttt{OnceLock} y
    el patrón \emph{lazy build}.
    \item Capítulo~\ref{cap:lftj}: SIMD intersection en LFTJ.
\end{itemize}

\noindent
Cuando aparezcan, vamos a referirnos de vuelta a este capítulo
sin re-derivar las definiciones.

\begin{ejercicios}
\begin{enumerate}
    \item \starred{} Una \emph{cache line} es de 64 B. Tu \texttt{struct}
    tiene 48 B útiles. Si lo declarás directo, dos instancias
    consecutivas comparten una cache line. ¿Bueno o malo? Pensalo
    bajo dos escenarios: (a) acceso secuencial; (b) acceso aleatorio
    desde threads distintos. (Pista: \emph{false sharing}.)

    \item Calculá: en una iteración secuencial de un
    \texttt{Vec<CompactRelation>} con 100M elementos, ¿cuántos cache
    line transfers son necesarios? Si la latencia efectiva es 100 ns
    por línea y el prefetch oculta el 90\%, ¿cuánto tarda en total?

    \item \starred{} ¿Por qué un \texttt{HashMap<K, V>} no tiene
    \emph{locality}? Diseñá una alternativa con
    \emph{open addressing} y \emph{linear probing} y argumentá si
    sería mejor o peor para una cache.

    \item Un mutex bajo contención cuesta $\sim$1\,$\mu$s por
    operación. Una atómica $\sim$10 ns. Si tu workload tiene 10K
    operaciones por segundo en 8 threads concurrentes, ¿cuál es la
    sobrecarga total con mutex? ¿Y con atómica? ¿Cambia la
    conclusión si las operaciones son \(99\%\) reads?

    \item Imaginá que \texttt{CompactRelation} pesara 24 bytes en vez
    de 32. Calculá: (a) cuántas aristas caben en L1 (32 KB); (b)
    cuánta memoria total ocuparían 100M aristas. ¿Por qué crees que
    el motor eligió 32 y no 24?

    \item Mapeás un archivo de 100~GB con mmap en una máquina de
    16~GB de RAM. ¿Funciona? ¿Qué pasa si el matcher hace un
    \emph{random access} sobre todo el archivo? ¿Cómo cambia la
    respuesta si el acceso es secuencial?
\end{enumerate}
\end{ejercicios}
